\documentclass[11pt]{article}
\usepackage[french]{babel}
\usepackage[utf8]{inputenc} 
\usepackage[T1]{fontenc}  
\usepackage{fullpage}
\usepackage{listings}
\usepackage{verbatim}
\usepackage{amssymb,amsmath}
\usepackage{url}
\usepackage{fullpage}


%************************************
%\topmargin -1.5cm
\textheight 25cm
\textwidth 18cm
\oddsidemargin -1.0cm

\newenvironment{note}{
\expandafter\comment
}{
\expandafter\endcomment
}

\newcommand{\correction}{
\renewenvironment{note}{
~\hspace*{-0.1cm} \hrulefill\\ \textbf{Notes.}~\\
}
{
~\hspace*{-0.1cm} \hrulefill
}
}
%\newcommand{\note}[1]{~\hspace*{-0.1cm} \hrulefill\\ \textbf{Notes.} #1 ~\\ ~\hspace*{-0.1cm} \hrulefill}
%\newcommand{\note}[1]{}
%************************************

\newtheorem{exercice}{Exercice}[section]
\newenvironment{sol}[1]{\par\noindent {\center \large \bf Proposition de solution pour l'exercice #1 \\~\\} \rm}{~\\ $\fbox{}$\bigskip} 

\title{Fondements informatiques I : Initiation au  Python\\
TD 2}
\begin{document}
\date{}
\maketitle
%\correction


\section{Révision Python}


%Exercice sur le printf et le scanf : formatage des chaînes

%Exercice d'appel de fonction -> pas fait en cours

%TD :
%expression logique : adresse ip IPV4 forme : x1.x2.x3.x4 entiers compris entre 0 et 255 (récupérés)
%obtenir adresse du réseau, adresse de broadcast à partir du masque de sous réseau


\bigskip

\begin{exercice}
Que font les programmes suivants?
\begin{enumerate}
\item 
\begin{verbatim}
L=[0 for i in range(5)]
LL=[]
for i in range(5):
    LL.append(L)

LL[0][1]=1
print(LL)    
\end{verbatim}

\item 
\begin{verbatim}
N=[]
for i in range(5):
    L=[0 for i in range(5)]
    N.append(L)

N[0][1]=1
print(N)  
\end{verbatim}

\item 
\begin{verbatim}
L=[0 for i in range(5)]
S=[L for i in range(5)]

S[0][1]=1
print(S)  
\end{verbatim}
\end{enumerate}
\end{exercice}

\section{Algorithmique et programmation en Python}


\begin{exercice}
Une file est une structure de données basée sur le principe que les premiers éléments ajoutés dans la file seront les premiers à en être retirés.
Dans cette exercice on utilisera une file d'entiers pour la gestion d'un stock de produits dans un entrepôt. 
La file contient des tuples indiquant à chaque fois le nombre de produits ainsi que leur jour d'arrivée:

\begin{verbatim}
[(100, 23), {159,22), ..., (45,2),(456,1),(100,1)]
\end{verbatim}
Le gestionnaire doit connaître à tout moment le nombre de produits disponibles. 

\begin{itemize}
\item[1.] Écrire une procédure  qui prend en argument la file, un entier correspondant au nombre de produits qui arrivent à l'entrepôt et un entier correspondant au stock courant. 
La procédure renverra la file modifiée et le stock courant. 
\item[2.]  Écrire une procédure qui prend en argument la file, un entier correspondant au nombre de produits qui doivent quitter l'entrepôt, et qui met à jour la file après avoir fait envoyer les produits 
conservés le plus longtemps. 
\item[3.] Donner le code d'un programme en Python pour les procédure des questions 1. et 2. Testez-le chez vous. 
 \end{itemize}



\end{exercice}




\begin{exercice}

Une pile est une structure de données basée sur le principe que le dernier élément ajouté dans la pile est le premier à en être retiré.
Nous l'utiliserons pour gérer une pile de livres à ranger sur une étagère. Chaque livre est représenté par un entier correspondant à son nombre de pages.

\vspace{0.3 cm}
\begin{itemize}
\item[1.] \'Ecrire le pseudo-code d'une procédure  \texttt{ajoute\_livre} qui prend en argument une liste et un entier \texttt{nb\_pages}, et qui ajoute le livre dans la liste (en haut de la pile, donc en fin de liste).
\item[2.] \'Ecrire le pseudo-code  d'une procédure \texttt{ranger\_livre} prend en argument la liste qui donne la pile, enlève le dernier élément ajouté (le livre en haut de la pile),
et renvoie le nombre de pages de ce livre.
\item[3.] \'Ecrire le pseudo-code d'un algorithme \texttt{ranger\_pile\_livres} qui prend en argument  la pile de livres, traite tous les livres dans la pile jusqu’à ce qu’elle soit vide,
et renvoie le total de pages de tous les livres rangés. 
\item[4.] Donner le code d'un programme en Python permettant de ranger une pile de livre. 
Testez (chez vous) votre programme après avoir  ajouté les livres suivants dans la pile (dans cet ordre) :


\vspace{0.5 cm}

\begin{verbatim}
150
200
120
80
220
\end{verbatim}


\end{itemize}

\end{exercice}




\begin{exercice}
\begin{itemize}
\item[1.]  Donner le pseudo-code d'un algorithme  qui construit  une matrice  \(5 \times 5\) contenant tous les nombres de 1 à 25 dans cet ordre: 

\[
\begin{matrix}
1 & 2 & 3 & 4 & 5 \\
6 & 7 & 8 & 9 & 10 \\
11 & 12 & 13 & 14 & 15 \\
16 & 17 & 18 & 19 & 20\\
21 & 22 & 23 & 24 & 25\\
\end{matrix}
\]


\item[2.]  Donner le pseudo-code d'un algorithme qui prend en entrée une matrice $5 \times 5$ et  construit une matrice dans laquelle: 
\begin{itemize}
\item la première ligne est décalée d'un cran vers la droite, et ceci de de manière circulaire,
\item la deuxième ligne est décalée de deux crans vers la droite, de manière circulaire,
\item la troisième est décalée de trois crans vers la droite, de manière circulaire, 
\item la quatrième est décalée de quatre crans vers la droite. de manière circulaire. 
\item la dernière ligne reste inchangée. 
\end{itemize}

Voici la sortie de l'algorithme pour la matrice de la question 1. 

\[
\begin{matrix}
 2 & 3 & 4 & 5 & 1 \\
 8 & 9 & 10  & 6 & 7\\
 14 & 15  & 11 & 12 & 13 \\
 20 & 16 & 17 & 18 & 19 \\
21 & 22 & 23 & 24 & 25\\
\end{matrix}
\]


\end{itemize}
\item Coder en Python les algorithmes des questions 1. et 2. 

\end{exercice}



\end{document}
